\documentclass[11pt,a4paper]{article}
\usepackage[margin=25mm]{geometry}
\usepackage{amsmath,amssymb,amsthm,booktabs,hyperref}
\setlength{\emergencystretch}{2em}
\newtheorem{theorem}{Theorem}\newtheorem{corollary}[theorem]{Corollary}
\newcommand{\Sym}{\operatorname{Sym}}
\title{Stability and finite starting cases for symmetric string C-groups}
\author{Human Compression project}\date{}
\begin{document}\maketitle
\begin{abstract}
We recover the principal results of Cameron, Fernandes and Leemans,
\emph{The number of string C-groups of high rank}, in its arXiv version
of 11 January 2023. The rank and degree correspondence is proved using
pure perfect cuts, actual subgroup intersections and an explicit
contraction. A distinguished pair of pure cuts reduces the six finite
starting counts to rooted markings on at most twelve points. The
resulting independent enumeration gives the original constants
$1,1,7,9,35,48$. The complete proofs of the two acquired components and
their supporting theory accompany this assembly; numerical constants
are not borrowed from the original table.
\end{abstract}

A string C-marking is an ordered list of nonidentity involutions
$\rho_0,\ldots,\rho_{r-1}$ satisfying distant commutation and the
intersection property
\[
 \langle\rho_i:i\in I\rangle\cap\langle\rho_i:i\in J\rangle
       =\langle\rho_i:i\in I\cap J\rangle
\]
for all label sets $I,J$. Marked-group isomorphisms preserve the ordered
marks. Reversal of their order is duality. Let $\Sigma^\kappa(n)$
be the markings of $\Sym(n)$ of rank $n-\kappa$, modulo marked-group
isomorphism and duality. This is the convention of the source's
introduction and Section 2.

\begin{theorem}[The principal stability statement]\label{stability}
For every $\kappa\ge1$ and $n\ge2\kappa+3$, there is a bijection
\[
                 \Sigma^\kappa(n)\longrightarrow
                 \Sigma^\kappa(n+1).
\]
It is obtained by subdividing an internal pure perfect cut. Its inverse
contracts the first internal pure pair. Both maps are well defined on
the stated isomorphism and duality classes. Thus every class at a later
degree is obtained from a class at the boundary degree $2\kappa+3$.
\end{theorem}
\begin{proof}
Apply the complete acquired component \emph{Pure cuts and the stable
counting correspondence}, Theorem 7. We record its exact interfaces to
make the application explicit. A pure perfect cut has mark $(a\ b)$,
and its preceding and following intervals are supported and transitive
on the two omission orbits containing $a,b$. For a fractured symmetric
marking with $2r>n+2$, the acquired sufficient-split induction and the
full leaf-kernel criterion give an internal pure perfect cut.
Fracture is automatic in the required range for $n\ge7$. Degrees five
and six use the ordinary simplex representative, with the exceptional
outer action of $S_6$ identified only at the level of abstract marked
groups.

The acquired proof shows that all pure perfect cuts form a single
consecutive path, with every internal path vertex fixed by the other
marks. Subdividing a pure edge inserts one point and one mark.
Actual one-bridge subgroup intersections prove that contracting an
adjacent pure pair preserves IP and generates the full symmetric group
on one fewer point. No quotient-IP inference is used. Exact omission
orbit transport shows that the two operations are inverse at the least
internal pure cut. The density inequality makes that cut begin an
internal pure pair in the target. Subdivision at any internal position
of the pure path gives an ordered-conjugate marking; this also pays
compatibility with reversal, including self-dual classes.

For the present parameters, $r=n-\kappa$ and
$2r\ge n+3$, so all domain hypotheses hold. In the target,
$N=n+1$, $R=r+1$, and $2R\ge N+4$, precisely the range used for
contraction. The maps therefore have the displayed domains and codomains.
Iteration gives the stated constructive reduction.
\end{proof}

\begin{theorem}[The six finite starting counts]\label{counts}
Writing $c_\kappa=|\Sigma^\kappa(2\kappa+3)|$, one has
\[
                         (c_1,\ldots,c_6)=(1,1,7,9,35,48).
\]
These are the values of $|\Sigma^\kappa(n)|$ for every
$n\ge2\kappa+3$ and $1\le\kappa\le6$.
\end{theorem}
\begin{proof}
The complete acquired component \emph{The finite starting cases from
balanced pure cuts}, Theorem 1, gives a unique consecutive pure pair
whose prefix surpluses are zero and one at the boundary $n=2r-3$.
Removing its middle vertex leaves two rooted transitive fractured
C-markings of degrees $2p,2q$ and ranks $p,q$, with $p+q=r-2$.
Each admits a pure leaf at its root. Conversely the exact leaf criteria,
the actual bridge intersection identity and pure subdivision reconstruct
the boundary marking. Ordered point conjugacy gives precisely rooted
point conjugacy on the two wings, and reversal exchanges the wings.

Theorem 3 of that component computes the rooted wing numbers by a
complete literal involution search, giving
$(w_1,\ldots,w_6)=(1,1,6,3,11,19)$. All first involutions are considered;
the search is over actual permutation tuples modulo the root stabilizer,
including degree-six actions without identifying them by outer group
automorphisms. The completeness proof, all representatives, terminal
receipts, independent labelled controls and actual IP checks are retained.
The unordered-pair formula consequently gives
\[
 c_\kappa=\sum_{\substack{p<q\\p+q=\kappa+1}}w_pw_q+
 \begin{cases}\binom{w_h+1}{2},&\kappa+1=2h,\\0,&\text{otherwise}.
 \end{cases}
\]
This yields respectively $1,1,6+1,3+6,11+3+21,19+11+18$.
Theorem~\ref{stability} transfers these boundary values to all degrees
in the stated ranges.
\end{proof}

\paragraph{Source contract and attribution.}
Theorem~\ref{stability} recovers source Theorems 1.1 and 9.10 and the
constructive classification reduction stated after Theorem 1.1.
Theorem~\ref{counts} recovers the six constants in the abstract,
introduction and the relevant boundary entries of Table 1. These were
already results of the original authors; no new resolution of an open
conjecture is claimed. Other entries of the historical table are not
silently added to this principal contract.

The original rank and degree operation, the rank classifications and
the established C-group criteria retain their original attribution.
The new proof uses backward propagation of full leaf kernels, a single
interval of pure cuts, explicit actual contraction and a balanced pair
for the finite count. The supporting recovered components distinguish
these deductions from reused proofs and computational inputs. The
prior accepted corrections to the unrestricted converse in source
Proposition 2.4, the indicated case of Proposition 2.5, and the unrestricted
extension claim in Theorem 6.4 remain in force. The proof above uses
only their exact repaired interfaces; it does not restore the false
unqualified statements. Those auxiliary counterexamples do not refute
the principal stability theorem.

The source edition is arXiv:2212.12723v2, 11 January 2023. A later
edition and the OEIS entry contain an additional numerical term; it is
outside this fixed edition's six-term contract and is not inferred here.
The separately registered classification paper L042 keeps its own
principal contract and original count.
\end{document}
